% Auto-generated by src/analysis/latex_tables.py -- do not edit by hand.
% LaTeX preamble required: \usepackage{booktabs, threeparttable, siunitx}
% Citations used: \citep{Galloway2007}.
\begin{table}[htbp]
\centering
\small
\begin{threeparttable}
\caption{Descriptive statistics by Catholic share}
\label{tab:descriptive_statistics}
\begin{tabular}{l S[table-format=3.2] S[table-format=3.2]}
\toprule
\textbf{Variable} & {Low-Cath} & {High-Cath} \\
 & {Mean} & {Mean} \\
\midrule
Crude Birth Rate (CBR)   & 38.52 & 39.77 \\
Legitimate Birth Rate   & 35.07 & 37.69 \\
Illegitimacy ratio (\%)   & 8.94 & 5.09 \\
General Marriage Rate (GMR)   & 12.77 & 12.34 \\
General Marital Fertility Rate (GMFR)   & 269.87 & 309.21 \\
Share Catholic (\%)   & 9.30 & 84.02 \\
\midrule
Counties (N) & {262} & {130} \\
Observations ($N \times T$) & \multicolumn{2}{c}{10,783 (Total)} \\
\bottomrule
\end{tabular}
\begin{tablenotes}
\footnotesize
\item \textit{Note:} Low-Catholic counties are defined as $\le 50\%$ Catholic in the 1871 census; High-Catholic as $> 50\%$. Means are pooled across the full 1862--1890 panel. The Crude Birth Rate and Legitimate Birth Rate are per 1{,}000 \emph{mid-year} inhabitants; the illegitimacy ratio is illegitimate births as a percentage of all births; the General Marriage Rate (GMR) is marriages per 1{,}000 women aged 15--49; and the General Marital Fertility Rate (GMFR) is legitimate births per 1{,}000 married women aged 15--49. The married-women and women 15--49 denominators of GMR and GMFR are time-varying, piecewise-linearly interpolated between Galloway's STA1871, AGE1882, and AGE1890 anchors. Temporal dynamics for each variable are shown in the time-series figures of the descriptive-evidence section.
\item \textit{Source:} Author's calculations from the Galloway Prussia Database \citep{Galloway2007}; mid-year population constructed via linear interpolation between consecutive December censuses.
\end{tablenotes}
\end{threeparttable}
\end{table}
